\documentclass[11pt]{article}
\usepackage[a4paper,margin=2.5cm]{geometry}
\usepackage{amsmath}
\usepackage{booktabs}
\usepackage{parskip}

\title{Moving-Average Trading Signal}
\author{}
\date{}

\begin{document}
\maketitle

\section*{Idea}
A moving-average signal is a simple rule-based classifier for a price time series. It transforms historical prices into a binary decision:
\[
    s_t =
    \begin{cases}
    1, & \text{invest / hold the stock},\\
    0, & \text{do not invest / stay out.}
    \end{cases}
\]

The signal compares a short-term moving average with a long-term moving average. The short-term average represents the recent price level, while the long-term average represents the broader baseline.

\section*{Moving Average Formula}
For a stock price series $P_t$, the moving average with window length $w$ is
\[
    MA_t^{(w)} = \frac{1}{w} \sum_{i=0}^{w-1} P_{t-i}.
\]
This is a trailing average, meaning that only the current and past prices are used.

\section*{Signal Rule}
The moving-average crossover signal is defined as
\[
    s_t =
    \begin{cases}
    1, & \text{if } MA_t^{(short)} > MA_t^{(long)},\\
    0, & \text{otherwise.}
    \end{cases}
\]

In words: if the recent average price is above the longer-term average price, the stock is classified as being in a positive trend.

\section*{Small Example}
Assume the prices over five days are:
\[
100,\; 101,\; 102,\; 105,\; 108.
\]
Using a short window of 3 days and a long window of 5 days, the moving averages on day 5 are
\[
    MA_5^{(3)} = \frac{102 + 105 + 108}{3} = 105,
\]
\[
    MA_5^{(5)} = \frac{100 + 101 + 102 + 105 + 108}{5} = 103.2.
\]
Since
\[
    105 > 103.2,
\]
the signal is
\[
    s_5 = 1.
\]
So the rule would decide to invest on day 5.

\section*{Data Science Interpretation}
The moving average itself is a feature-engineering step. The full signal is the decision rule that compares two engineered features:
\[
    \text{short-term trend feature} > \text{long-term baseline feature}.
\]
Therefore, the signal behaves like a simple rule-based binary classifier for time-series data.


\section{Volatility-Filtered Time-Series Momentum Signal}

The signal is motivated by Moskowitz, Ooi, and Pedersen (2012), who document
time-series momentum: an asset's own past return can contain information about
its future return. Their original strategy uses the sign of the asset's past
return to determine the trading direction and scales the position size by
inverse ex-ante volatility.

Let \(P_t\) denote the adjusted close price at time \(t\). We define the
lookback return over \(k\) trading days as

\[
R_{t-k,t} = \frac{P_t}{P_{t-k}} - 1.
\]

Daily returns are defined as

\[
r_t = \frac{P_t}{P_{t-1}} - 1.
\]

Recent volatility over a window of \(w\) trading days is estimated by

\[
\sigma_t =
\sqrt{
\frac{1}{w}
\sum_{i=0}^{w-1}
(r_{t-i} - \bar{r}_t)^2
}.
\]

Because our assessment framework uses binary long-only signals, we adapt the
original long/short and volatility-scaled framework into a 0/1 rule:

\[
s_t =
\begin{cases}
1, & \text{if } R_{t-k,t} > 0 \text{ and } \sigma_t < c, \\
0, & \text{otherwise}.
\end{cases}
\]

Thus, the strategy invests only when the asset has positive recent momentum and
recent volatility is below the chosen threshold \(c\). This volatility filter is
not identical to the inverse-volatility position sizing in Moskowitz, Ooi, and
Pedersen (2012), but it is a simplified adaptation that fits the binary signal
structure of our assessment notebook.


\section{Trading-Range Breakout Signal}

The trading-range breakout signal is motivated by Brock, Lakonishok, and
LeBaron (1992), who study simple technical trading rules on the Dow Jones
Index. Besides moving-average rules, they explicitly analyze trading-range
break rules. In their description, a buy signal is generated when the price
breaks above a resistance level, where the resistance level is defined as a
local maximum. They also test versions based on past 50, 150, and 200 trading
days and versions with and without a one percent band \cite[p.~1736]{brock1992}.

For our binary long-only assessment framework, we adapt this idea into an
invest-or-stay-out rule. Let $P_t$ denote the adjusted close price at time $t$.
For a breakout window of length $n$, the recent resistance level is defined as

\[
H_t^{(n)} = \max\{P_{t-n}, P_{t-n+1}, \ldots, P_{t-1}\}.
\]

The current price $P_t$ is intentionally not included in this maximum. This
avoids look-ahead bias and makes the rule compare today's price only with
information that was already available before day $t$.

The breakout strength can be written as

\[
B_t^{(n)} = \frac{P_t}{H_t^{(n)}} - 1.
\]

A positive value of $B_t^{(n)}$ means that today's price is above the highest
price observed during the previous $n$ trading days. The binary trading signal
is then defined as

\[
s_t =
\begin{cases}
1, & \text{if } P_t > H_t^{(n)}, \\
0, & \text{otherwise.}
\end{cases}
\]

In words, the strategy invests only when the stock reaches a new recent high.
This makes the breakout signal different from the moving-average signal. The
moving-average signal compares the price level with an average baseline, while
the breakout signal compares the price level with an extreme historical value,
namely the recent maximum. Therefore, the breakout rule is stricter: a stock can
be above its moving average without breaking above its previous high.

\subsection*{Economic Reasoning}

The economic intuition is that a breakout above the recent maximum may indicate
strong buying pressure and trend continuation. In technical-analysis language,
the recent maximum is interpreted as a resistance level: if the price rises above
this level, the previous selling pressure has been overcome. Brock, Lakonishok,
and LeBaron (1992) describe this as a trading-range break-out rule and motivate
it by the idea that prices breaking above previous peaks can generate buy
signals \cite[p.~1736]{brock1992}.

For our project, the signal is important because it provides a second
trend-following rule that is not based on a rolling average. From a data science
perspective, it is a transparent rule-based classifier using an engineered
extreme-value feature. The assessment requires economic reasoning, reasonable
statistics, and convincing empirical evidence for the chosen signals and their
parameters. This signal is appropriate because its parameter $n$ has a clear
interpretation as the number of past trading days used to define the resistance
level.

\subsection*{Small Example}

Assume the following price series:

\[
100,\; 110,\; 90,\; 101.
\]

Using a breakout window of $n=3$, the resistance level on day 4 is

\[
H_4^{(3)} = \max\{100, 110, 90\} = 110.
\]

The breakout strength is

\[
B_4^{(3)} = \frac{101}{110} - 1 = -0.0818.
\]

Since

\[
101 < 110,
\]

the trading-range breakout signal is

\[
s_4 = 0.
\]

Thus, the strategy does not invest on day 4. This example also shows why the
breakout rule is stricter than a moving-average rule. The three-day average of
$100$, $110$, and $90$ is $100$, so the price $101$ is above the moving average.
However, it is still below the recent maximum of $110$, so there is no breakout.

\subsection*{Parameter Choice}

Following the original paper, natural candidate windows are $n=50$, $n=150$,
and $n=200$ trading days \cite[p.~1736]{brock1992}. In our empirical notebook,
we can additionally test a shorter window such as $n=20$, which approximately
represents one trading month. The parameters can then be compared in-sample and
evaluated out-of-sample using the same performance statistics and graphs as for
the other trading signals.

\section{Short-Term Reversal Signal}

The short-term reversal signal is motivated by the empirical literature on
predictable return reversals. The main economic idea is different from the
trend-following signals above. Moving-average, volatility-filtered momentum,
and breakout signals invest after price strength. The reversal signal instead
invests after short-term price weakness.

Jegadeesh (1990) provides evidence that individual stock returns are predictable
from past returns. In the abstract and empirical results, he reports highly
significant negative first-order serial correlation in monthly stock returns
\cite[p.~881]{jegadeesh1990}. This means that, at the one-month horizon, high
past returns tend to be followed by lower future returns and low past returns
tend to be followed by higher future returns. In his portfolio test labelled
strategy $S1$, securities are ranked in ascending order based on their one-month
lagged returns \cite[pp.~887--888]{jegadeesh1990}. This is the direct reversal
logic: recent losers are placed in the highest predicted-return portfolio.

Lehmann (1990) studies the same mechanism at a shorter weekly horizon. He states
that one-week winners and losers experience sizeable reversals in the following
week \cite[p.~1]{lehmann1990}. His empirical strategy gives negative weights to
recent winners and positive weights to recent losers, so that the portfolio earns
profits if short-term returns reverse \cite[p.~3]{lehmann1990}. In the original
cross-sectional formulation, the portfolio weight is proportional to the negative
of the stock's return relative to the equal-weighted market return:

\[
W_{i,t-k}
=
-\left(R_{i,t-k} - \bar{R}_{t-k}\right),
\]

where $R_{i,t-k}$ is the return of stock $i$ in the previous period and
$\bar{R}_{t-k}$ is the equal-weighted average return across the stock universe
\cite[p.~4]{lehmann1990}. If a stock was a winner relative to the average, then
$R_{i,t-k} - \bar{R}_{t-k} > 0$ and the weight becomes negative. If a stock was a
loser relative to the average, then $R_{i,t-k} - \bar{R}_{t-k} < 0$ and the
weight becomes positive. Thus, the original paper implements a long-short
cross-sectional reversal strategy.

Our assessment notebook uses single-asset, long-only binary signals. Therefore,
we adapt the same reversal mechanism into an invest-or-stay-out rule. Let
$P_t$ denote the adjusted close price and let $k$ be the lookback window. The
recent return is

\[
R_{t,k} = \frac{P_t}{P_{t-k}} - 1.
\]

A negative value of $R_{t,k}$ means that the stock has lost value over the last
$k$ trading days. The entry rule is defined as

\[
e_t =
\begin{cases}
1, & \text{if } R_{t,k} < \theta, \\
0, & \text{otherwise},
\end{cases}
\]

where $\theta < 0$ is the reversal threshold. For example, $\theta=-0.05$ means
that the signal is triggered only after a loss of more than five percent over the
lookback window.

To make the signal closer to the predictive logic of the papers, the strategy
uses a fixed holding period $h$. If an entry signal occurs at time $t$, the
strategy remains invested for the next $h$ signal observations:

\[
s_{t+j}=1
\quad \text{for} \quad
j = 0,1,\ldots,h-1
\quad \text{if} \quad e_t = 1.
\]

Otherwise, the strategy stays out of the stock:

\[
s_t = 0.
\]

In the Python implementation, overlapping entry signals are allowed. If a new
entry occurs while the strategy is already invested, the signal simply remains
at one for the relevant holding-period interval.

\subsection*{Economic Reasoning}

The economic interpretation is that large negative short-term returns may partly
reflect temporary overreaction, liquidity pressure, or order-flow imbalances
rather than a permanent decline in fundamental value. If the price decline is
partly temporary, the following period may contain a partial rebound. This is why
Lehmann's original strategy buys recent losers and shorts recent winners, and why
Jegadeesh's one-month reversal strategy ranks securities with low lagged returns
more favorably.

For our project, the signal is useful because it provides a genuinely different
economic mechanism from the previous signals. The previous rules are mainly
trend-following. This rule is contrarian: it buys after weakness rather than
after strength. From a data science perspective, the engineered feature is a
lagged return, and the classifier activates when this feature falls below a
negative threshold.

\subsection*{Small Example}

Assume the following prices over six trading days:

\[
100,\;98,\;97,\;96,\;95,\;94.
\]

Let the lookback window be $k=5$ and the threshold be $\theta=-0.05$. On day 6,
the five-day return is

\[
R_{6,5} = \frac{94}{100} - 1 = -0.06.
\]

Since

\[
-0.06 < -0.05,
\]

the entry signal is

\[
e_6 = 1.
\]

With a holding period of $h=5$, the strategy remains invested for five signal
observations after the reversal event. In words: the stock lost more than five
percent over one trading week, so the strategy classifies the stock as strongly
oversold and invests in expectation of a short-term correction.

\subsection*{Parameter Choice}

Natural candidate lookback windows are $k=5$, $k=10$, and $k=20$, corresponding
approximately to one week, two weeks, and one trading month. Natural thresholds
are $\theta=-0.03$, $\theta=-0.05$, $\theta=-0.08$, and $\theta=-0.10$. A simple
baseline choice is

\[
k=5, \qquad \theta=-0.05, \qquad h=5.
\]

This means: if the stock loses more than five percent over the previous trading
week, invest for the next trading week. The parameters should be calibrated
in-sample and then evaluated out-of-sample using the same performance statistics
and graphs as the other signals.

\begin{thebibliography}{9}

\bibitem{brock1992}
Brock, W., Lakonishok, J., and LeBaron, B. (1992).
\textit{Simple Technical Trading Rules and the Stochastic Properties of Stock Returns}.
The Journal of Finance, 47(5), 1731--1764.

\bibitem{moskowitz2012}
Moskowitz, T. J., Ooi, Y. H., and Pedersen, L. H. (2012).
\textit{Time Series Momentum}.
Journal of Financial Economics, 104(2), 228--250.


\bibitem{jegadeesh1990}
Jegadeesh, N. (1990).
\textit{Evidence of Predictable Behavior of Security Returns}.
The Journal of Finance, 45(3), 881--898.

\bibitem{lehmann1990}
Lehmann, B. N. (1990).
\textit{Fads, Martingales, and Market Efficiency}.
The Quarterly Journal of Economics, 105(1), 1--28.

\end{thebibliography}


\end{document}
